\noindent The project adopts an interdisciplinary approach combining computer science, archaeology, cultural heritage studies, and digital humanities.
\noindent Technological development is closely integrated with heritage expertise, ensuring both scientific accuracy and user accessibility. This approach enables effective collaboration across domains and enhances the societal relevance of the project outcomes.

\vspace{1em}
\noindent This interdisciplinary synergy is manifested through the following pillars:
\begin{itemize}
    \item \textbf{Technology and Publishing:} Arts et Métiers (CAPLAB) provides cutting-edge XR tools, while Springer Nature ensures that the resulting digital heritage trend reaches the widest possible academic and public audience through open science and open access mechanisms.
    \item \textbf{Mediterranean Heritage Experts:} Sapienza Università di Roma and the National and Kapodistrian University of Athens bridge classical Greco-Roman archaeology, urban conservation, and climate risk assessments with modern digital documentation methods.
    \item \textbf{European Regional Synergy:} ELTE BTK (Hungary) brings AI-based processing and digital humanities expertise (Danube-AI), while Sorbonne Université (France) coordinates data collection on medieval fortifications and coastal erosion, ensuring a truly pan-European and cross-disciplinary framework.
\end{itemize}